%% Detailed measurements (moved out of the deployed-instantiation section for length).
\section{Detailed measurements}
\label{app:meas}

This appendix expands the headline of \Cref{sec:results} with the full prover,
verifier, and proof sweeps, the per-phase decomposition, the beyond-RAM
behaviour, and the head-to-head comparison with \binius{}. All figures use the
setup of \Cref{sec:results} (one Apple~M4, matched rate $1/4$, CPU unless noted).

\subsection{Comparison to Binius64 at a matched commitment rate}
\label{sec:impl-binius}

Binius64 is the closest relative of this construction: another
characteristic-$2$ SHA-256 prover that refuses the prime-field witness
blow-up, and the source of the per-step full-adder carry identity and the
one-$\AND$ Boolean forms used in \Cref{sec:arith}. The two systems commit
differently --- Binius64 over a Reed--Solomon code with a recursive FRI
proximity test, the deployed instantiation here over the linear $\Ftwo$-RAA
code with a flat Brakedown opening (\Cref{sec:commitment}) --- but they
share one tunable dial, the code rate, which trades commit work against
query count, hence proof size. Binius64 runs at rate $1/2$ by default; the
deployed instantiation runs at rate $1/4$ (\Cref{sec:impl-params}). A
\emph{prover} comparison is therefore only meaningful at a fixed rate.
\Cref{tab:binius-rate} reports both systems on a matched workload of
$\approx\!963$ compressions ($\nu=16$): at rate $1/4$, and --- for Binius64
--- at its native rate $1/2$. All figures are CPU-only medians on one
machine, each system at its deployed soundness query count (here $t=987$
column openings of \Cref{sec:impl-params}; Binius64's test-query count is
fixed automatically from its $96$-bit target). The comparison is CPU-only
on both sides; enabling the Metal GPU commit of \Cref{sec:commitment}
offloads the commit phase to the GPU but, at this scale, shows no net
prover-latency win over the CPU commit (\Cref{tab:scaling-gpu}).

\begin{table}[h]
\centering
\renewcommand{\arraystretch}{1.2}
\begin{tabular}{l r r r}
\toprule
System (rate) & Prove & Verify & Proof \\
\midrule
Zinc\texttt{+} (rate $1/4$) & $\sim\!93$ ms & $\sim\!12$ ms & $\sim\!1.1$ MiB \\
Binius64 (rate $1/4$) & $\sim\!104$ ms & $\sim\!15$ ms & $\sim\!249$ KiB \\
Binius64 (rate $1/2$, default) & $\sim\!95$ ms & $\sim\!14$ ms & $\sim\!304$ KiB \\
\bottomrule
\end{tabular}
\caption{SHA-256 at $\nu=16$ ($\approx\!963$ compressions), CPU-only,
representative medians. At a matched rate $1/4$ the $\Ftwo[X]$ prover is
$\sim\!1.1\times$ faster than Binius64, and is no slower than Binius64 at
its cheaper native rate $1/2$. Raising Binius64 to rate $1/4$ adds $\approx\!8$ ms to its prover ---
a $2\times$ larger codeword to commit and fold --- while shrinking its
proof, since its $\AND$/shift sum-checks are rate-independent and only the
commitment scales with the rate.}
\label{tab:binius-rate}
\end{table}

\begin{remark}[Reading the comparison]
\label{rem:binius-read}
At equal rate the linear-time RAA encoder and the base-field Hadamard
discharge (\Cref{rem:cost-structure}) make the $\Ftwo[X]$ prover the faster
of the two; Binius64's recursive FRI in turn keeps its proof several times
smaller, and the two verifiers are comparable at this scale. Two
asymmetries temper the prover figure: Binius64 is fully sound and proves a
length-padded preimage, whereas the deployed adder Hadamards here are
supplied trusted (\Cref{rem:adder-trusted}) --- closing that gap is
expected to be near prover-neutral but is not yet paid; and the Binius64
build measured here carries a parallelized verifier. As elsewhere, the
ratios are steadier than the absolute milliseconds.
\end{remark}

\subsection{Scaling: prover, verifier, and proof}
\label{sec:impl-scaling}\label{sec:res-scaling}

\Cref{tab:scaling} sweeps both systems across trace sizes at the matched rate
$1/4$ --- Binius64 by compression count, Zinc\texttt{+} by trace height $\nu$
($\approx(2^\nu-4)/68$ compressions). The Zinc\texttt{+} figures are CPU-only
(no GPU commit) on a cooled machine; Binius64 is its own example at
\texttt{--log-inv-rate 2}. Two divergences stand out. The Zinc\texttt{+}
prover stays ahead and the gap \emph{widens} --- from $\sim\!1.1\times$ at
$10^3$ compressions to $\sim\!2\times$ at $10^4$. And the verifiers split
qualitatively: Binius64's grows with $N$ (its public-input check is $O(N)$,
reaching $\sim\!0.5$ s at $8000$ compressions), while Zinc\texttt{+}'s ---
dominated by the $t$ column openings --- grows like $\sqrt N$ and stays in the
tens of milliseconds. Binius64's example also walls above $\sim\!8000$
compressions: its $O(N)$ circuit \emph{build} exceeds one minute, where
Zinc\texttt{+} still proves $15{,}420$ compressions in $\approx 1.2$ s. The
Metal GPU commit configuration tracks the same curve without a net latency
win at these scales (\Cref{tab:scaling-gpu}).

\begin{table}[h]
\centering
\renewcommand{\arraystretch}{1.2}
\begin{tabular}{r r r r r}
\toprule
 & \multicolumn{2}{c}{Binius64 (rate $1/4$)} & \multicolumn{2}{c}{Zinc\texttt{+} (CPU)} \\
\cmidrule(lr){2-3}\cmidrule(lr){4-5}
$\approx$compressions & prove & verify & prove & verify \\
\midrule
$\sim\!60$    & $33.7$ ms & $3.1$ ms  & $10.1$ ms & $4.8$ ms \\
$\sim\!250$   & $41.6$ ms & $4.6$ ms  & $23.3$ ms & $9.3$ ms \\
$\sim\!1000$  & $84.1$ ms & $14.0$ ms & $75.0$ ms & $17.9$ ms \\
$\sim\!2000$  & $155$ ms  & $27.5$ ms & $148$ ms  & $19.8$ ms \\
$\sim\!8000$  & $1.22$ s  & $493$ ms  & $581$ ms  & $45$ ms \\
$\sim\!15{,}400$ & \multicolumn{2}{c}{(build $>1$ min)} & $1.19$ s & $83$ ms \\
\bottomrule
\end{tabular}
\caption{SHA-256 scaling at rate $1/4$, CPU, representative single runs.
\emph{Proof sizes} over the same range: Binius64 $123\!\to\!299$ KiB
($16\!\to\!8192$ compressions; recursive FRI grows slowly); Zinc\texttt{+}
grows as $O(\sqrt N)$ --- $0.55/2.88/10.3$ MiB at $\approx\!7/964/15{,}420$
compressions (the balanced bit-sliced opening of \Cref{tab:zinc-scaling}, in
which the combined row and the $t=987$ column openings are within $\sqrt2$ of
each other at the optimal shape). Verify figures
are single-run and indicative (the $\nu{=}16$ median is the $\approx\!12$ ms
of \Cref{tab:binius-rate}); the $O(N)$-vs-$\sqrt N$ verifier split is the
robust signal. Binius64 cannot go below $16$ compressions (its example's fixed
$1024$-byte message).}
\label{tab:scaling}
\end{table}

\Cref{tab:zinc-scaling} gives the full deployed Zinc\texttt{+} sweep, CPU-only,
at the proof-size-optimal matrix shape ($\mathrm{num\_rows}\approx\sqrt N$;
\Cref{sec:commitment}). Proof size grows as $O(\sqrt N)$ --- the balanced shape
trades the combined row against the $t$ column openings, so each is within
$\sqrt 2$ of the other and the total is twice their geometric mean. It is
therefore \emph{not} flat: from $0.55$ MiB at one block to $\approx 40$ MiB at
$2^{24}$ rows. (A fixed, small $\mathrm{num\_rows}$ would instead make the
proof $O(N)$; balancing the shape is what keeps it $O(\sqrt N)$.) The prover and
verifier scale near-linearly and $\approx\sqrt N$ respectively until $\nu\ge 23$,
where the working set exceeds the $16$ GB host and both become memory-bound.

\begin{table}[h]
\centering
\renewcommand{\arraystretch}{1.15}
\begin{tabular}{r r r r r r}
\toprule
$\nu$ & $\approx$comp. & $\mathrm{num\_rows}$ & prove & verify & proof \\
\midrule
$9$  & $7$         & $2$   & $2.2$ ms  & $1.2$ ms  & $0.55$ MiB \\
$14$ & $241$       & $16$  & $13.6$ ms & $3.8$ ms  & $1.62$ MiB \\
$16$ & $964$       & $32$  & $46.6$ ms & $8.7$ ms  & $2.88$ MiB \\
$19$ & $7{,}710$   & $64$  & $354$ ms  & $38$ ms   & $7.41$ MiB \\
$20$ & $15{,}420$  & $128$ & $667$ ms  & $64$ ms   & $10.33$ MiB \\
$21$ & $30{,}841$  & $128$ & $1.35$ s  & $117$ ms  & $14.36$ MiB \\
$22$ & $61{,}683$  & $256$ & $2.72$ s  & $224$ ms  & $20.21$ MiB \\
$23$ & $123{,}366$ & $256$ & $7.0$ s$^\dagger$  & $383$ ms$^\dagger$ & $28.23$ MiB \\
$24$ & $246{,}733$ & $512$ & $25.1$ s$^\dagger$ & $1.00$ s$^\dagger$ & $39.93$ MiB \\
\bottomrule
\end{tabular}
\caption{Deployed Zinc\texttt{+} SHA-256 scaling: CPU-only, rate $1/4$, at the
proof-size-optimal shape $\mathrm{num\_rows}=\arg\min\bigl(\mathrm{row\_len}\cdot
s_{\mathrm{row}} + t\,b\,\mathrm{num\_rows}\cdot s_{\mathrm{cell}}\bigr)\approx
\sqrt{N\,s_{\mathrm{row}}/(t\,b\,s_{\mathrm{cell}})}$ ($s_{\mathrm{row}}=512$ B,
$s_{\mathrm{cell}}=8$ B, $t=987$, $b=6$ committed columns). Apple~M-series
medians (thermal-dependent; ratios are stabler than absolutes). $^\dagger$At
$\nu\ge 23$ the working set exceeds the $16$ GB RAM, so the prove/verify figures
are memory-bound (super-linear); proof size is unaffected.}
\label{tab:zinc-scaling}
\end{table}

\Cref{tab:scaling-gpu} reports the same Zinc\texttt{+} prover with the Metal
GPU commit enabled, swept across the deployed trace heights.

\begin{table}[h]
\centering
\renewcommand{\arraystretch}{1.2}
\begin{tabular}{r r r}
\toprule
$\nu$ & $\approx$compressions & prove (GPU commit) \\
\midrule
$9$  & $7$        & $9.8$ ms \\
$12$ & $60$       & $19.2$ ms \\
$14$ & $240$      & $38.8$ ms \\
$16$ & $963$      & $107$ ms \\
$18$ & $3855$     & $403$ ms \\
$20$ & $15{,}420$ & $1.71$ s \\
\bottomrule
\end{tabular}
\caption{Zinc\texttt{+} prover with the Metal GPU commit enabled
(\Cref{sec:commitment}), three-run medians via the discharge timing harness.
The curve is linear in $N$ (each $+1$ in $\nu$ doubles the compression count
and $\approx$doubles the prover), matching the shape of the CPU prover of
\Cref{tab:scaling}. The absolutes track the \emph{warm} CPU median of
\Cref{tab:binius-rate} ($\nu{=}16$: $107$ vs $93$ ms) and run above the
\emph{cooled} CPU sweep of \Cref{tab:scaling} --- so once thermal state is
matched, offloading the commit to the GPU yields no net end-to-end latency
win: its value is throughput and freeing the CPU, not reduced prove latency.
(This harness rebuilds the $\mathrm{GF}(2^8)$ lookup tables per verify call,
inflating its verify above the deployed $\approx\!12$ ms; the prover figure
is unaffected.)}
\label{tab:scaling-gpu}
\end{table}

\paragraph{Beyond RAM.} \Cref{tab:scaling-bigmem} pushes the Zinc\texttt{+}
prover past the point where the trace exceeds the host's $16$\,GB of memory.
Two regimes are visible. While the trace fits in RAM (through $\nu{=}20$,
$\approx\!13$\,GB), the prover is \emph{cleanly linear} --- $2.0\times$ per
doubling of $N$. From $\nu{=}21$ ($\approx\!26$\,GB) the working set exceeds
physical memory and the prover turns super-linear ($2.3$--$3.1\times$ per
doubling) as memory compression and SSD swap absorb the overflow. This knee
measures \emph{this host's memory wall}, not the algorithm, whose compute
stays $O(N)$; on a machine with adequate RAM the linear regime simply extends.
The salient point for the comparison is that Zinc\texttt{+} still
\emph{completes} $123{,}361$ compressions --- a $\approx\!104$\,GB working set,
entirely via swap --- where Binius64's $O(N)$ circuit \emph{build} already walls
at $\approx\!16{,}000$ (\Cref{tab:scaling}); the all-$\Ftwo[X]$ prover stays
tractable an order of magnitude further in $N$, even on memory-constrained
hardware.

\begin{table}[h]
\centering
\renewcommand{\arraystretch}{1.2}
\begin{tabular}{r r r r l}
\toprule
$\nu$ & $\approx$comp & prove & verify & working set \\
\midrule
$19$ & $7{,}710$   & $0.59$ s  & $45$ ms  & fits \\
$20$ & $15{,}420$  & $1.18$ s  & $96$ ms  & $\sim\!13$ GB, fits \\
\midrule
$21$ & $30{,}840$  & $3.51$ s  & $126$ ms & $\sim\!26$ GB (swap) \\
$22$ & $61{,}680$  & $8.02$ s  & $313$ ms & $\sim\!52$ GB (swap) \\
$23$ & $123{,}361$ & $25.2$ s  & $663$ ms & $\sim\!104$ GB (swap) \\
\bottomrule
\end{tabular}
\caption{Zinc\texttt{+} prover beyond RAM (one Apple~M4, $16$\,GB; CPU, merged
discharge, $t=987$, same-session warm medians). Above the rule the trace fits
memory and the prover is linear ($2.0\times$ per doubling); below it the working
set exceeds $16$\,GB and the prover turns super-linear ($2.3$--$3.1\times$),
from memory compression and SSD swap --- a property of the host, not the $O(N)$
algorithm. Zinc\texttt{+} nonetheless completes all the way to $123{,}361$
compressions, where Binius64's circuit build has long since walled
(\Cref{tab:scaling}). The verifier continues its $\sqrt N$ growth (tens to
hundreds of ms) across the whole range.}
\label{tab:scaling-bigmem}
\end{table}

\subsection{Per-phase scaling}
\label{sec:impl-perphase}\label{sec:res-phases}

\Cref{tab:perphase} decomposes the Zinc\texttt{+} prover by phase across trace
heights. Every phase scales near-linearly --- each grows $\approx\!4\times$ per
$4\times$ compressions, and the total tracks $N^{0.97}$ --- so the prover is
linear-time with no super-linear term, the property that lets it run past the
size where a non-uniform circuit's $O(N)$ build stalls. The ends differ: the
\emph{commit} (rate-$1/4$ codeword $+$ Merkle) scales steepest and its share
grows from $7\%$ to $30\%$ --- it is the large-$N$ centre of gravity, and the
phase the GPU commit most relieves --- while the \emph{open} scales shallowest
(its $t=987$ column openings are a \emph{fixed} count), its share shrinking
from $17\%$ to $9\%$.

\begin{table}[h]
\centering
\renewcommand{\arraystretch}{1.2}
\begin{tabular}{l r r r c}
\toprule
Phase & $\nu{=}10$ & $\nu{=}16$ & $\nu{=}20$ & per $4\times$ \\
      & ($15$c) & ($963$c) & ($15{,}420$c) & comp \\
\midrule
Commit (encode $+$ Merkle)          & $0.4$ & $18.0$ & $337.9$ & $4.3\times$ \\
Linear PIOP (IC, $\psia$, sumcheck) & $1.9$ & $17.8$ & $255.3$ & $3.8\times$ \\
Hadamard discharge                  & $1.2$ & $19.0$ & $287.6$ & $3.9\times$ \\
Multipoint eval                     & $1.1$ & $11.2$ & $146.0$ & $3.6\times$ \\
Open                                & $0.9$ & $9.6$  & $102.7$ & $3.3\times$ \\
\midrule
Total                               & $5.3$ & $75.6$ & $1129$  & $3.9\times$ \\
\bottomrule
\end{tabular}
\caption{Zinc\texttt{+} merged prover by phase (ms), CPU, $t=987$, across trace
height $\nu$. Near-linear throughout ($4\times$ per $4\times$ compressions is
exactly linear); commit scales steepest, open shallowest. Single profile runs;
the per-phase \emph{shares} are stable across thermal state.}
\label{tab:perphase}
\end{table}

\medskip
\noindent\textbf{Binius64, for contrast.} The same sweep on Binius64
(\Cref{tab:perphase-b64}) scales differently. Its prover is
\emph{super-linear} and worsening --- $\sim\!4.6\times$ per $4\times$
compressions around $10^3$, steepening past $6\times$ by $8000$ (versus
Zinc\texttt{+}'s flat $\sim\!3.9\times$) --- because the BitAnd and Shift
reductions build $\mathrm{GF}(2^{128})$ tensors whose working set spills
cache. The Shift reduction dominates at small $N$ but its share falls
($62\%\!\to\!32\%$); the BitAnd reduction rises to dominate
($14\%\!\to\!41\%$, its witness build alone a quarter of the prover) and
drives the super-linear growth, while the FRI commit grows with the codeword.

\begin{table}[h]
\centering
\renewcommand{\arraystretch}{1.2}
\begin{tabular}{l r r r r}
\toprule
Phase & $64$c & $1024$c & $4096$c & $8192$c \\
\midrule
Commit (FRI)        & $1.8$  & $15.9$ & $86.3$  & $144.7$ \\
BitAnd reduction    & $4.0$  & $32.7$ & $219.4$ & $557.1$ \\
Shift reduction     & $16.5$ & $51.9$ & $171.2$ & $434.7$ \\
PCS opening         & $3.9$  & $13.4$ & $45.9$  & $212.7$ \\
\midrule
Total prove         & $27.3$ & $116$  & $532$   & $1370$ \\
\bottomrule
\end{tabular}
\caption{Binius64 prover by phase (ms), rate $1/4$. The BitAnd reduction
overtakes the Shift reduction at large $N$ and drives super-linear growth.
IntMul is a $\le\!1$ ms no-op for SHA.}
\label{tab:perphase-b64}
\end{table}

\medskip
\noindent\textbf{Preprocessing data.} The starkest structural difference is
not in any timed phase but in the \emph{preprocessing} --- the circuit
description the system holds (and that the verifier in effect must read).
Binius64 is a non-uniform word circuit, so its constraint system plus the
shift/wiring key-collection grow $O(N)$: measured at rate $1/4$,
$3.7 / 14.4 / 54.7 / 218 / 873$ MB at $16 / 64 / 256 / 1024 / 4096$
compressions (almost exactly $\times4$ per $\times4$), and the
$\ge\!16{,}384$-compression build exceeds one minute. Zinc\texttt{+} is a
\emph{uniform} AIR: the constraint description is the same for any trace
height, so its preprocessing is $O(1)$ --- a fixed UAIR specification plus the
RAA code's permutation seeds, kilobyte-scale and independent of the
compression count. This $O(N)$-versus-$O(1)$ preprocessing gap is the root of
both Binius64's $O(N)$ verifier and its circuit-build wall.

